\documentclass[11pt]{article}
\usepackage{mathblatt}
\begin{document}
% Kopfzeile Seite 1 ohne Kennung; die Rückseite (Lösungen, für den Lehrer) bekommt sie nach \begleitteil
\blattkopf{Mathematik}{Basisaufgaben}
\einheitenkopf[][Zettel 11]{Basisaufgaben · Zettel 11}
\zweigzeile{je Aufgabe 1 Punkt · Lösungen auf der Rückseite}
\par\vspace{2pt}\noindent Name \underline{\hspace{7cm}}\hfill Datum \underline{\hspace{3cm}}\hfill Punkte \underline{\hspace{1.2cm}} / 18\par\vspace{6pt}
% \small: volle Seite, kurze Aufgaben zweispaltig (v0.6)
\small

\noindent\begin{minipage}[t]{0.485\linewidth}\raggedright
% 1: Bruchteil einer Größe berechnen – brueche-dezimalzahlen-basis-k2-v4 (Original 2018-OS-B1a, ausgedacht: keine Marke)
\begin{aufgabe}{$\frac{5}{6}$ von $1{,}2$ kg – wie viel Kilogramm? \leerfeld[kg]}
\end{aufgabe}
\end{minipage}\hfill\begin{minipage}[t]{0.485\linewidth}\raggedright
% 2: Zahlen in verschiedenen Darstellungen vergleichen – brueche-dezimalzahlen-basis-k4-v9 (Original 2023-OS-B1f, ausgedacht: keine Marke)
\begin{aufgabe}{Welcher Wert ist der größte: $0{,}1$; $0{,}1^2$; $10\,\%$; $\frac{1}{8}$? \leerfeld}
\end{aufgabe}
\end{minipage}\par

\noindent\begin{minipage}[t]{0.485\linewidth}\raggedright
% 3: Termwert berechnen – rationale-zahlen-basis-k1-v7 (Original 2026-FOR-B1g, ausgedacht: keine Marke)
\begin{aufgabe}{Berechne den Wert von $4 \cdot (x - 6)$ für $x = -2$. \leerfeld}
\end{aufgabe}
\end{minipage}\hfill\begin{minipage}[t]{0.485\linewidth}\raggedright
% 4: Exponent einer Potenz bestimmen – potenzen-wurzeln-basis-k1-v4 (Original 2024-OS-B1g, ausgedacht: keine Marke)
\begin{aufgabe}{$2^x = 64$ – welche Zahl ist $x$? x = \leerfeld}
\end{aufgabe}
\end{minipage}\par

\noindent\begin{minipage}[t]{0.485\linewidth}\raggedright
% 5: Prozentwert berechnen – prozentrechnung-basis-k4-v7 (Original 2026-FOR-B1a, ausgedacht: keine Marke)
\begin{aufgabe}{$70\,\%$ von $50\,€$ \leerfeld[€]}
\end{aufgabe}
\end{minipage}\hfill\begin{minipage}[t]{0.485\linewidth}\raggedright
% 6: Größen vergleichen – einheiten-basis-k2-v1 (Original 2026-FOR-B1d, ausgedacht: keine Marke)
\begin{aufgabe}{Vergleiche $0{,}8$ km und $80$ m. Setze das richtige Zeichen ein: $<$, $=$ oder $>$. $0{,}8$ km \leerfeld $80$ m}
\end{aufgabe}
\end{minipage}\par

\noindent\begin{minipage}[t]{0.485\linewidth}\raggedright
% 7: Antiproportionale Zuordnung Dreisatz – zuordnungen-basis-k1-v3 (Original 2014-GYM-B1c, ausgedacht: keine Marke)
\begin{aufgabe}{6 Pumpen leeren ein Becken in 10 Stunden. Wie lange brauchen 5 Pumpen? \leerfeld[h]}
\end{aufgabe}
\end{minipage}\hfill\begin{minipage}[t]{0.485\linewidth}\raggedright
% 8: Lineare Gleichung lösen – lineare-gleichungen-basis-k1-v6 (Original 2024-OS-B1d, ausgedacht: keine Marke)
\begin{aufgabe}{Löse die Gleichung $3 \cdot (x - 2) = 12$. x = \leerfeld}
\end{aufgabe}
\end{minipage}\par

% 9: Lösung durch Einsetzen prüfen – quadratische-gleichungen-basis-k1-v5 (Original 2025-OS-B1h, ausgedacht: keine Marke)
\begin{aufgabe}{Welcher Wert erfüllt $x \cdot (x - 2) = 3$? \\ \kreuz{$x = 1$} \kreuz{$x = 3$} \kreuz{$x = -3$} \kreuz{$x = -2$}}
\end{aufgabe}

% 10: Grundwert berechnen – prozentrechnung-basis-k1-v3 (Original 2025-OS-B1a, ausgedacht: keine Marke)
\begin{aufgabe}{Beim Kauf eines Kopfhörers spart Ida $12\,€$, das sind $10\,\%$ des alten Preises. Wie hoch war der alte Preis? \leerfeld[€]}
\end{aufgabe}

% 11: Figur nach Spiegelung benennen – symmetrie-abbildungen-basis-k1-v3 (Original 2018-OS-B1h, ausgedacht: keine Marke)
\begin{aufgabe}{Ein gleichschenkliges Dreieck liegt mit der Basis auf der Geraden $g$ und wird an $g$ gespiegelt. Welches Viereck entsteht?}
\end{aufgabe}

% 12: Dreiecksungleichung anwenden – winkel-dreiecke-basis-k1-v3 (Original 2020-OS-B1i, ausgedacht: keine Marke)
\begin{aufgabe}{Ein dreieckiges Segel ABC hat die Seite $c = 4$ m. Was gilt für die Summe der beiden anderen Seiten $a$ und $b$? \\ \kreuz{$a + b > 4$ m} \kreuz{$a + b = 4$ m} \kreuz{$a + b < 4$ m}}
\end{aufgabe}

% 13: Eigenschaft einer Figur zuordnen – winkel-dreiecke-basis-k5-v1 (Original 2026-FOR-B1c, ausgedacht: keine Marke)
\begin{aufgabe}{Kreuze die richtige Ergänzung an: „In jedem Rechteck …“\\ \kreuz{… sind die Diagonalen gleich lang.}\\ \kreuz{… stehen die Diagonalen senkrecht aufeinander.}\\ \kreuz{… sind alle Seiten gleich lang.}}
\end{aufgabe}

% 14: Arithmetisches Mittel berechnen – daten-basis-k1-v1 (Original 2021-OS-B1f, ausgedacht: keine Marke)
\begin{aufgabe}{Eine Radtour dauert drei Tage. Gefahren werden $42$; $35$; $49$ km. Berechne das arithmetische Mittel (Strecke pro Tag). \leerfeld[km]}
\end{aufgabe}

% 15: Fehlenden Wert aus Mittelwert bestimmen – daten-basis-k2-v1 (Original 2022-OS-B1d, ausgedacht: keine Marke)
\begin{aufgabe}{Vier Schüler haben im Test diese Punkte: $18$; $22$; \underline{\hspace{1cm}}; $20$. Der Durchschnitt beträgt $21$. Ergänze den fehlenden Wert. \leerfeld Punkte}
\end{aufgabe}

% 16: Ergebnismenge aufzählen – wahrscheinlichkeit-basis-k1-v1 (Original 2017-OS-B1f, ausgedacht: keine Marke)
\begin{aufgabe}{Man wirft eine Münze (Zahl oder Wappen) und einen Spielwürfel. Wie viele verschiedene Ergebnisse gibt es? \leerfeld Ergebnisse}
\end{aufgabe}

% 17: Bruchteil einer Fläche bestimmen – brueche-dezimalzahlen-basis-k1-v10 (Original 2026-FOR-B1b, ausgedacht: keine Marke)
\begin{aufgabe}{Bei welcher Figur ist genau $\frac{3}{5}$ grau? \\ \kreuz{P} \kreuz{Q} \kreuz{R} \kreuz{S}}

P \bruchrechteck[2.5]{6}{10} \quad Q \kreissektor[0.8]{200}{} \quad R \bruchkreis[0.8]{3}{6} \quad S \bruchrechteck[2.5]{3}{4}
\end{aufgabe}

% 18: Term zu Figur angeben – flaechen-basis-k4-v1 (Original 2025-OS-B1i, ausgedacht: keine Marke)
\begin{aufgabe}{\begin{minipage}[t]{0.6\linewidth}Welche Formel beschreibt den Umfang $u$? \\ \kreuz{$u = x \cdot y$} \kreuz{$u = 2 \cdot x + 2 \cdot y$} \kreuz{$u = x + y$} \kreuz{$u = 4 \cdot x$}\end{minipage}\hfill\begin{minipage}[t]{0.37\linewidth}\vspace{0pt}
\rechteck[punkte=,seiten={x,y,,}]{3}{1.8}
\end{minipage}}
\end{aufgabe}

\begleitteil
\blattkopf{Basisaufgaben · Zettel 11}{BAS-Z11 · Lösungen}
\erg{1}{$1{,}2 : 6 \cdot 5 = 1$ kg \hfill {\footnotesize\textit{Bruchteil einer Größe berechnen · 2018}}}
\erg{2}{$\frac{1}{8} = 0{,}125$; die anderen sind $0{,}1$, $0{,}01$ und $0{,}1$ \hfill {\footnotesize\textit{Zahlen in verschiedenen Darstellungen vergleichen · 2023}}}
\erg{3}{$4 \cdot (-2 - 6) = 4 \cdot (-8) = -32$ \hfill {\footnotesize\textit{Termwert berechnen · 2026}}}
\erg{4}{$2, 4, 8, 16, 32, 64$: sechs Faktoren, also $x = 6$ \hfill {\footnotesize\textit{Exponent einer Potenz bestimmen · 2024}}}
\erg{5}{$0{,}7 \cdot 50 = 35\,€$ \hfill {\footnotesize\textit{Prozentwert berechnen · 2026}}}
\erg{6}{$0{,}8$ km $= 800$ m, also $0{,}8$ km $>$ $80$ m \hfill {\footnotesize\textit{Größen vergleichen · 2026}}}
\erg{7}{$6 \cdot 10 = 60$ Pumpenstunden; $60 : 5 = 12$ h \hfill {\footnotesize\textit{Antiproportionale Zuordnung Dreisatz · 2014}}}
\erg{8}{$x - 2 = 4$, also $x = 6$ \hfill {\footnotesize\textit{Lineare Gleichung lösen · 2024}}}
\erg{9}{$x = 3$, denn $3 \cdot (3 - 2) = 3 \cdot 1 = 3$ \hfill {\footnotesize\textit{Lösung durch Einsetzen prüfen · 2025}}}
\erg{10}{$12 \cdot 10 = 120\,€$ (nicht $10\,\%$ von $12\,€$) \hfill {\footnotesize\textit{Grundwert berechnen · 2025}}}
\erg{11}{Raute \hfill {\footnotesize\textit{Figur nach Spiegelung benennen · 2018}}}
\erg{12}{$a + b > 4$ m \hfill {\footnotesize\textit{Dreiecksungleichung anwenden · 2020}}}
\erg{13}{… sind die Diagonalen gleich lang. \hfill {\footnotesize\textit{Eigenschaft einer Figur zuordnen · 2026}}}
\erg{14}{$126 : 3 = 42$ km \hfill {\footnotesize\textit{Arithmetisches Mittel berechnen · 2021}}}
\erg{15}{$4 \cdot 21 = 84$; $84 - 60 = 24$ \hfill {\footnotesize\textit{Fehlenden Wert aus Mittelwert bestimmen · 2022}}}
\erg{16}{$2 \cdot 6 = 12$ (Z1 bis Z6 und W1 bis W6) \hfill {\footnotesize\textit{Ergebnismenge aufzählen · 2017}}}
\erg{17}{P, denn $\frac{6}{10} = \frac{3}{5}$ \hfill {\footnotesize\textit{Bruchteil einer Fläche bestimmen · 2026}}}
\erg{18}{$u = 2 \cdot x + 2 \cdot y$ \hfill {\footnotesize\textit{Term zu Figur angeben · 2025}}}
\end{document}
